\documentclass[11pt, a4paper]{article}

% --- PREAMBLE ---
\usepackage[a4paper, top=2.5cm, bottom=2.5cm, left=2cm, right=2cm]{geometry}
\usepackage{fontspec}
\usepackage{amsmath}
\usepackage{amssymb}
\usepackage{graphicx}
\usepackage{hyperref}

% --- BIBLIOGRAPHY ---
\usepackage[numbers, sort&compress]{natbib}

% --- LANGUAGE & FONTS ---
% We use babel for proper hyphenation
\usepackage[english]{babel}

% Font Configuration (Safe for Overleaf)
% If 'Noto Sans' fails, comment these two lines out to use default fonts.
\setmainfont{Noto Sans} 
\setsansfont{Noto Sans}

% --- METADATA ---
\title{Physics-Informed Residual Learning for Adaptive Control\\
\large Phase 1: Theory \& Modeling}
\author{Basyirin Amsyar bin Basri}
\date{February 2026}

\begin{document}

\maketitle

\section{Motivation and Related Work}

Standard neural networks cannot efficiently learn the high-frequency transients
of stick-slip friction without violating real-time inference constraints.
\citet{rahaman2019spectralbiasneuralnetworks} mathematically demonstrated the
spectral bias of ReLU networks, where the Fourier components $f_{\theta}(k)$
decay proportional to $k^{-(d+1)}$, causing the network to prioritize
low-frequency structure during training.
Because stick-slip transitions are severe discontinuities, they contain
infinite high-frequency components; a standard MLP therefore acts as an
implicit low-pass filter, blurring the stiction jump.
To force a network to capture these sharp transients from scratch, the
parameter count would need to scale beyond the capabilities of a
600\,MHz Cortex-M7 while maintaining a \textless10\,ms control loop.
To overcome this spectral bias without blowing up computational overhead, we
must manually inject the high-frequency jump via a physics prior.

While high-fidelity dynamic friction models exist, their inherent numerical
stiffness makes them unsuitable for sub-10\,ms embedded control.
The LuGre model introduces an internal friction state governed by
\begin{equation}
    \frac{dz}{dt} = v - \frac{\lvert v \rvert}{g(v)}\,z,
    \label{eq:lugre}
\end{equation}
requiring numerical integration (e.g., Forward Euler) on
hardware~\citep{376053}.
However, at near-zero velocities, the high bristle stiffness
($\sigma_0 \approx 10^5$\,N/m) renders the ODE mathematically
``stiff,'' requiring extremely small time steps to prevent divergence.
On a Teensy\,4.1, these extra CPU cycles for state-tracking and stiff-solver
logic destroy the deterministic 100\,Hz loop.
By utilizing the memoryless, algebraic Stribeck model as a prior and leaving a
lightweight residual learner to handle the complex transients, we achieve
accuracy without the computational penalty of an ODE solver.

Current residual control architectures have not successfully bridged the
sub-10\,ms embedded gap, remaining tethered to high-latency general-purpose
compute stacks.
While recent efforts have applied PINNs to hardware control, they either
bottleneck control loops to sub-10\,Hz by offloading computation to
Python/JAX environments incompatible with bare-metal
deployment~\citep{liu2023physicsinformedneuralnetworksmodel}, or achieve
high-frequency estimation only by inheriting the non-deterministic scheduling
jitter of x86 Linux architectures tethered to desktop
compute~\citep{10971218}.
Relying on these computing stacks introduces operating system scheduling jitter
and overhead that is incompatible with bare-metal constraints; even when
achieving a 1\,kHz loop~\citep{10971218}, running a PINN on a tethered x86
laptop fundamentally fails the embedded sim-to-real criterion, as it cannot be
deployed autonomously on low-power, edge-native microcontrollers like the
Cortex-M7.
This paper directly addresses this gap by demonstrating a PIRL architecture
implemented in native C++ on a Teensy\,4.1, achieving deterministic
\textgreater100\,Hz control without reliance on tethered desktop compute.

\section{System Dynamics and Problem Formulation}

\subsection{Nominal Dynamics}
We consider a single-degree-of-freedom (1-DoF) rotational actuator. The fundamental equation of motion governing the system is given by:

\begin{equation}
    J \ddot{q} + F_f(\dot{q}) = \tau_{u} - \tau_{ext}
    \label{eq:motion}
\end{equation}

where $q, \dot{q}, \ddot{q}$ represent the angular position, velocity, and acceleration, respectively. $J$ is the moment of inertia, $\tau_{u}$ is the control input (motor torque), and $\tau_{ext}$ represents external disturbances. The term $F_f(\dot{q})$ denotes the nonlinear friction torque, which is the primary source of modeling uncertainty in low-cost actuators.

\subsection{The Stribeck Friction Model}
Standard viscous damping models ($F_f = B\dot{q}$) are insufficient for precision control at low velocities due to the presence of Coulomb friction and the Stribeck effect (stick-slip phenomenon). We adopt the continuous Stribeck model to capture these dynamics:

\begin{equation}
    F_f(\dot{q}) = \text{sgn}(\dot{q}) \left[ F_c + (F_s - F_c) e^{-|\dot{q}/v_s|^\delta} \right] + \sigma \dot{q}
    \label{eq:stribeck}
\end{equation}

The parameters are defined as:
\begin{itemize}
    \item $F_c$: Coulomb friction torque (dynamic sliding friction).
    \item $F_s$: Static friction torque (breakaway torque).
    \item $v_s$: Stribeck velocity (threshold where stiction transitions to sliding).
    \item $\delta$: Shape factor (typically $\delta = 2$ for Gaussian decay).
    \item $\sigma$: Viscous friction coefficient.
\end{itemize}

This nonlinearity creates a "dead zone" at zero velocity, making pure PID control ineffective (steady-state error or oscillation).

\subsection{Proposed Control Architecture}
To compensate for the uncertainty in Eq. (\ref{eq:stribeck}), we propose a composite control law. We assume the friction consists of a nominal linear component (known) and a nonlinear residual (unknown):

\begin{equation}
    F_f(\dot{q}) = \underbrace{\hat{\sigma} \dot{q}}_{\text{Nominal}} + \underbrace{\Delta(\dot{q})}_{\text{Residual}}
\end{equation}

The control input $\tau_u$ is defined as:

\begin{equation}
    \tau_u = \tau_{PID} + \tau_{PINN}(\dot{q})
\end{equation}

where:
\begin{itemize}
    \item $\tau_{PID}$: A standard Proportional-Integral-Derivative controller responsible for tracking the nominal error dynamics.
    \item $\tau_{PINN}$: A Physics-Informed Neural Network (PINN) estimator trained to predict the residual term $\Delta(\dot{q})$ (the "Rust").
\end{itemize}

% --- FIGURE PLACEHOLDER ---
\begin{figure}[h]
    \centering
    \includegraphics[width=0.8\textwidth]{stribeck_curve_week2.png}
    \caption{Visualization of the Stribeck Friction Curve showing the nonlinear transition from static ($F_s$) to Coulomb ($F_c$) friction.}
    \label{fig:stribeck}
\end{figure}


\bibliographystyle{plainnat}
\bibliography{references}

\end{document}